\chapter{Center and Spread}
\label{ch:center-spread}

A dashboard shows five request latencies from one server, in milliseconds:
\[
  20,\quad 22,\quad 25,\quad 28,\quad 35.
\]
Two questions come up immediately. What is a \emph{typical} latency? And how
\emph{different} are the requests from each other? The first question asks for a
measure of \textbf{center}; the second asks for a measure of \textbf{spread}. Every
summary in this chapter answers one of those two questions, and each answer has a
precise sense in which it is the ``right'' one. Chapter~1 drew the whole shape of a
data set as a histogram; this chapter squeezes that shape into two or three numbers
you can put on a dashboard.

Notation: the data are $x_1, x_2, \dots, x_n$. Sorted in increasing order they are
written $x_{(1)} \le x_{(2)} \le \dots \le x_{(n)}$. These are the \emph{order
statistics}: $x_{(1)}$ is the smallest value, $x_{(n)}$ the largest, and the bracket
in the subscript means ``position after sorting.''

\section{The mean}
\label{sec:center-spread-mean}

\begin{definition}[Sample mean]
The \textbf{mean} (arithmetic average) of $x_1,\dots,x_n$ is
\[
  \xbar = \frac{1}{n}\sum_{i=1}^n x_i .
\]
\end{definition}

For the latencies,
\[
  \xbar = \frac{20+22+25+28+35}{5} = \frac{130}{5} = 26 \text{ ms}.
\]
No request actually took 26 ms. The mean need not be a data value; it is a
summary, not a member.

\subsection{The mean is the balance point}

Put each observation as an equal weight on a weightless ruler. The ruler balances
at exactly one place, and that place is $\xbar$. The reason is one line of algebra.
Call $d_i = x_i - \xbar$ the \textbf{deviation} of observation $i$. Then
\[
  \sum_{i=1}^n d_i = \sum_{i=1}^n x_i - n\xbar = n\xbar - n\xbar = 0.
\]
Deviations to the left exactly cancel deviations to the right. For the latencies
the deviations are $-6,-4,-1,+2,+9$: the left side totals $-11$ and the right side
totals $+11$. In physics terms the total \emph{torque} (weight times signed distance
from the pivot, the quantity that makes a seesaw tip) about $\xbar$ is zero.

\begin{figure}[H]
\centering
\begin{tikzpicture}[x=0.55cm,y=1cm]
  \draw[thick] (18,0) -- (37,0);
  \foreach \t in {18,20,...,36} {
    \draw (\t,0) -- (\t,-0.12) node[below,font=\scriptsize] {\t};
  }
  \foreach \v in {20,22,25,28,35} {
    \fill[sourceblue] (\v,0.22) circle (0.2);
  }
  \fill[quotientorange] (26,-0.55) -- (25.45,-1.25) -- (26.55,-1.25) -- cycle;
  \node[quotientorange,anchor=west,font=\small] at (26.7,-0.95) {$\xbar=26$};
  \draw[-{Stealth},trapred] (26,0.75) -- (20,0.75) node[midway,above,font=\scriptsize] {$-6$};
  \draw[-{Stealth},trapred] (26,1.25) -- (22,1.25) node[midway,above,font=\scriptsize] {$-4$};
  \draw[-{Stealth},trapred] (26,1.75) -- (25,1.75) node[left,font=\scriptsize] {$-1$};
  \draw[-{Stealth},targetgreen] (26,1.25) -- (28,1.25) node[right,font=\scriptsize] {$+2$};
  \draw[-{Stealth},targetgreen] (26,0.75) -- (35,0.75) node[midway,above,font=\scriptsize] {$+9$};
  \draw[dashed,gray] (26,0) -- (26,2.1);
  \node[font=\scriptsize,anchor=west] at (29.5,1.9) {left: $-11$,\; right: $+11$};
  \node[font=\scriptsize] at (27.5,-1.85) {latency (ms)};
\end{tikzpicture}
\caption{The mean is where the ruler balances. Signed deviations from $\xbar=26$
sum to zero: $-6-4-1+2+9=0$.}
\label{fig:center-spread-balance}
\end{figure}

The balance picture explains the mean's main weakness. A weight far from the rest
has a long lever arm. Move the 35 ms request to 180 ms and the balance point moves
from 26 to $(20+22+25+28+180)/5 = 55$ ms, above four of the five values.

\begin{recoverybox}
Pause and think: add a sixth request that took exactly 26 ms. Where does the ruler
balance now? (It stays at 26. The new weight sits on the pivot, so it adds no torque.)
\end{recoverybox}

\section{The median, the mode, and the trimmed mean}
\label{sec:center-spread-median}

\begin{definition}[Median]
Sort the data. If $n$ is odd, the \textbf{median} is the middle value
$x_{((n+1)/2)}$. If $n$ is even, it is the average of the two middle values,
$\tfrac12\big(x_{(n/2)} + x_{(n/2+1)}\big)$.
\end{definition}

The five latencies have median $25$ (the third of five). Add a sixth request of
180 ms. Now $n=6$, the middle pair is $25$ and $28$, and the median is $26.5$.
The mean jumps from $26$ to $(130+180)/6 \approx 51.7$.

\begin{figure}[H]
\centering
\begin{tikzpicture}[x=0.065cm,y=1cm]
  \foreach \yy/\lab in {1.6/{$n=5$},0/{$n=6$}} {
    \draw[thick] (10,\yy) -- (190,\yy);
    \node[anchor=east,font=\small] at (8,\yy) {\lab};
  }
  \foreach \t in {20,40,...,180} {
    \draw (\t,0) -- (\t,-0.1) node[below,font=\scriptsize] {\t};
  }
  \foreach \v in {20,22,25,28,35} {
    \fill[sourceblue] (\v,1.75) circle (1.3pt);
    \fill[sourceblue] (\v,0.15) circle (1.3pt);
  }
  \fill[trapred] (180,0.15) circle (1.3pt);
  \draw[targetgreen,thick] (25,1.35) -- (25,2.15);
  \node[targetgreen,font=\scriptsize,anchor=south] at (21,2.15) {median 25};
  \draw[quotientorange,thick] (26,1.35) -- (26,2.15);
  \node[quotientorange,font=\scriptsize,anchor=south west] at (30,1.95) {mean 26};
  \draw[targetgreen,thick] (26.5,-0.25) -- (26.5,0.55);
  \node[targetgreen,font=\scriptsize,anchor=south] at (26.5,0.55) {median 26.5};
  \draw[quotientorange,thick] (51.7,-0.25) -- (51.7,0.55);
  \node[quotientorange,font=\scriptsize,anchor=west] at (53,0.45) {mean 51.7};
  \node[font=\scriptsize] at (100,-0.75) {latency (ms)};
\end{tikzpicture}
\caption{One slow request (red, 180 ms) drags the mean from 26 to 51.7 ms, past
five of the six values. The median barely moves, from 25 to 26.5.}
\label{fig:center-spread-median-outlier}
\end{figure}

The median ignores \emph{how far} the extreme values are. It only uses their
\emph{rank}. That is why ``p50 latency'' (the 50th percentile: half of requests are faster) on a
monitoring dashboard is the median: a
single pathological request cannot distort it.

\begin{definition}[Mode]
The \textbf{mode} is the most frequent value.
\end{definition}

The mode is the only center that works for categorical data. In a log of HTTP
status codes $200, 200, 404, 200, 500, 200, 404$ the mode is $200$. A mean of
status codes is meaningless. For continuous measurements every value is usually
distinct, so the mode is read off a histogram as the tallest bar (the ``modal
class''), and a distribution with two clear peaks is called \textbf{bimodal}.

\begin{definition}[Trimmed mean]
The \textbf{$\alpha$-trimmed mean} sorts the data, deletes the $\lfloor \alpha n\rfloor$
smallest and $\lfloor \alpha n \rfloor$ largest values, and averages what remains.
\end{definition}

The $0\%$ trimmed mean is the ordinary mean. Trimming just under $50\%$ from each
end leaves only the middle and gives the median. A trimmed mean sits between the
two: it uses the magnitudes of most of the data but discards the extremes.
Olympic diving and figure-skating scores use exactly this rule.

\section{Mean versus median: skew and outliers}
\label{sec:center-spread-skew}

\begin{example}[Salaries at a small company]
Ten employees earn, in thousands of dollars per year,
\[
  40,\ 45,\ 48,\ 50,\ 52,\ 55,\ 58,\ 60,\ 65,\ 900.
\]
The last salary is the founder's. The sum is $473 + 900 = 1373$, so
\[
  \xbar = 137.3, \qquad \Med = \frac{52+55}{2} = 53.5 .
\]
Nine of the ten people earn less than the mean. The $10\%$ trimmed mean drops one
value from each end (the 40 and the 900) and averages the remaining eight:
$433/8 = 54.125$. Without the founder the mean would be $473/9 \approx 52.6$ and
the median $52$.
\end{example}

\begin{figure}[H]
\centering
\begin{tikzpicture}
\begin{axis}[
  width=12.5cm, height=3.6cm,
  xmin=0, xmax=1000, ymin=-0.5, ymax=2.2,
  axis y line=none, axis x line=bottom,
  xlabel={salary (\$ thousands)},
  xtick={0,100,...,1000}, tick label style={font=\scriptsize},
  clip=false,
]
\addplot[only marks, mark=*, mark size=2pt, sourceblue]
  coordinates {(40,0.2)(45,0.45)(48,0.7)(50,0.95)(52,1.2)(55,0.2)(58,0.45)(60,0.7)(65,0.95)(900,0.2)};
\draw[targetgreen,thick] (axis cs:53.5,0) -- (axis cs:53.5,1.7);
\node[targetgreen,font=\scriptsize,anchor=south] at (axis cs:53.5,1.7) {median 53.5};
\draw[quotientorange,thick] (axis cs:137.3,0) -- (axis cs:137.3,1.7);
\node[quotientorange,font=\scriptsize,anchor=south west] at (axis cs:137.3,1.2) {mean 137.3};
\node[trapred,font=\scriptsize,anchor=south] at (axis cs:900,0.4) {founder};
\end{axis}
\end{tikzpicture}
\caption{Salaries at a ten-person company (points stacked to avoid overlap). One
large value puts the mean far to the right of where nine people actually are. The
median stays inside the cluster.}
\label{fig:center-spread-salary}
\end{figure}

Which number is ``correct''? Both, for different questions. If the company wants
to know its total payroll, the mean answers it: total $= n\xbar = 1373$. If a job
candidate wants to know what a typical employee earns, the median answers it. The
mean is the right summary for \emph{totals}; the median is the right summary for
\emph{a typical individual}.

\subsection{Skew moves the mean toward the tail}

A distribution is \textbf{right-skewed} (positively skewed) if it has a long tail to
the right: incomes, house prices, latencies, file sizes, city populations. The tail
values pull the mean right while the median stays with the bulk. The typical
ordering is
\[
  \text{mode} < \text{median} < \text{mean} \qquad\text{(right skew)},
\]
reversed for left skew (exam scores with a ceiling, age at death in rich countries).
For a symmetric distribution all three coincide.

\begin{recoverybox}
Pause and think: a news story reports that the mean house price in a city is
\$620{,}000 and the median is \$450{,}000. Which way is the distribution skewed, and
roughly what share of houses sell for less than \$620{,}000? (Right-skewed; more than
half, since the median is below the mean.)
\end{recoverybox}

\begin{figure}[H]
\centering
\begin{tikzpicture}
\begin{axis}[
  width=11.5cm, height=5.6cm,
  domain=0:8, samples=120,
  xmin=0, xmax=8, ymin=0, ymax=0.48,
  axis lines=left, xlabel={value}, ylabel={density},
  tick label style={font=\scriptsize}, clip=false,
]
\addplot[fill=histfill!50, draw=distblue, thick] {x*exp(-x)} \closedcycle;
\draw[structurepurple,thick] (axis cs:1,0) -- (axis cs:1,0.368);
\draw[targetgreen,thick] (axis cs:1.678,0) -- (axis cs:1.678,0.314);
\draw[quotientorange,thick] (axis cs:2,0) -- (axis cs:2,0.271);
\node[structurepurple,font=\scriptsize,anchor=south] at (axis cs:1,0.375) {mode 1};
\draw[targetgreen] (axis cs:1.678,0.30) -- (axis cs:3.2,0.36) node[anchor=west,font=\scriptsize] {median 1.68};
\draw[quotientorange] (axis cs:2,0.25) -- (axis cs:3.5,0.27) node[anchor=west,font=\scriptsize] {mean 2};
\node[font=\scriptsize,anchor=west] at (axis cs:4.3,0.09) {long right tail pulls the mean};
\end{axis}
\end{tikzpicture}
\caption{A right-skewed density curve, $f(x)=xe^{-x}$ (area under the curve between
two values is the share of data between them). The peak
(mode) is at 1, half the area lies left of 1.68 (median), and the balance point
(mean) is at 2. The tail drags the mean furthest.}
\label{fig:center-spread-skew}
\end{figure}

\begin{trapbox}
``The average household income is \$X'' does not say which center was used. If it
is a mean, then because incomes are right-skewed, well over half of households earn
less than \$X. When a distribution is skewed, ask which center was used. Government statistics on
income and house prices usually report the median for this reason.
\end{trapbox}

\subsection{Robustness and the breakdown point}

How much bad data can a statistic absorb before it becomes useless? Keep the nine
ordinary salaries and let the founder's salary $v$ grow. The median stays at
$53.5$ for every $v \ge 55$. The mean is $(473+v)/10$, which grows without bound.

\begin{figure}[H]
\centering
\begin{tikzpicture}
\begin{axis}[
  width=11cm, height=5.4cm,
  xmin=50, xmax=1000, ymin=0, ymax=160,
  axis lines=left,
  xlabel={founder's salary $v$ (\$ thousands)}, ylabel={summary value},
  legend pos=north west, legend style={font=\scriptsize},
  tick label style={font=\scriptsize},
]
\addplot[quotientorange, very thick, domain=65:1000] {(473+x)/10};
\addlegendentry{mean $=(473+v)/10$}
\addplot[targetgreen, very thick, domain=65:1000] {53.5};
\addlegendentry{median $=53.5$}
\addplot[structurepurple, very thick, dashed, domain=65:1000] {54.125};
\addlegendentry{10\% trimmed mean}
\end{axis}
\end{tikzpicture}
\caption{One contaminated value. The mean is a straight line in $v$ and can be made
as large as you like. The median ($53.5$) and the trimmed mean ($54.1$) are flat and
nearly on top of each other: no single value can move them far.}
\label{fig:center-spread-breakdown}
\end{figure}

\begin{definition}[Breakdown point]
The \textbf{breakdown point} of a statistic is the smallest fraction of the data
that, if replaced by arbitrary values, can make the statistic arbitrarily large.
\end{definition}

\begin{itemize}
  \item \textbf{Mean:} one bad value out of $n$ suffices, so the breakdown point is
  $1/n$, which tends to $0$.
  \item \textbf{Median:} you must corrupt half the data to drag the middle value
  away. Breakdown point $\approx 50\%$, the maximum possible.
  \item \textbf{$\alpha$-trimmed mean:} the trimming absorbs up to $\alpha n$ bad
  values at each end. Breakdown point $\approx \alpha$.
\end{itemize}

A statistic with a high breakdown point is called \textbf{robust}. Robustness has a
price: the median throws away the magnitudes of most observations, so for clean,
symmetric, bell-shaped data it is noisier than the mean: from sample to sample its
standard deviation is about $25\%$ larger (Chapter~6 shows this by simulation). Choose the mean when
the data are clean and you care about totals; choose the median when outliers or
heavy tails are plausible.

\section{Deviations, variance, and standard deviation}
\label{sec:center-spread-variance}

Two servers can share a mean latency of 26 ms while one is steady and the other
erratic. Spread measures how far observations typically sit from the center.

The natural starting point is the deviations $d_i = x_i - \xbar$. Their average is
useless: it is always $0$ (Section~\ref{sec:center-spread-mean}). Positive and
negative deviations cancel. Two fixes remove the signs: take absolute values, or
square. Squaring gives the variance.

\begin{definition}[Sample variance and standard deviation]
\[
  s^2 = \frac{1}{n-1}\sum_{i=1}^n (x_i - \xbar)^2, \qquad s = \sqrt{s^2}.
\]
$s^2$ is the \textbf{sample variance}; $s$ is the \textbf{sample standard
deviation} (SD). The divisor $n-1$ is explained in
Section~\ref{sec:center-spread-bessel}.
\end{definition}

\subsection{A worked table}

Build $s$ for the five latencies column by column.

\begin{center}
\begin{tabular}{c c r r}
\toprule
$i$ & $x_i$ (ms) & $d_i = x_i - 26$ & $d_i^2$ \\
\midrule
1 & 20 & $-6$ & 36 \\
2 & 22 & $-4$ & 16 \\
3 & 25 & $-1$ & 1 \\
4 & 28 & $+2$ & 4 \\
5 & 35 & $+9$ & 81 \\
\midrule
sum & 130 & $0$ & 138 \\
\bottomrule
\end{tabular}
\end{center}

\begin{enumerate}
  \item Mean: $\xbar = 130/5 = 26$.
  \item Deviations: subtract 26. Check that they sum to $0$.
  \item Square each deviation and add: $\sum d_i^2 = 138$ (the \textbf{sum of
  squares}).
  \item Divide by $n-1 = 4$: $s^2 = 138/4 = 34.5 \text{ ms}^2$.
  \item Take the square root: $s = \sqrt{34.5} \approx 5.87$ ms.
\end{enumerate}

\begin{figure}[H]
\centering
\begin{tikzpicture}[x=0.55cm,y=0.55cm]
  \draw[thick] (17,0) -- (37,0);
  \foreach \t in {18,20,...,36} \draw (\t,0) -- (\t,-0.2) node[below,font=\scriptsize] {\t};
  \draw[dashed,quotientorange] (26,-0.2) -- (26,9.6) node[above,font=\scriptsize] {$\xbar=26$};
  \fill[trapred!25,draw=trapred] (20,0) rectangle (26,6);
  \fill[trapred!45,draw=trapred] (22,0) rectangle (26,4);
  \fill[trapred!70,draw=trapred] (25,0) rectangle (26,1);
  \fill[targetgreen!30,draw=targetgreen] (26,0) rectangle (35,9);
  \fill[targetgreen!60,draw=targetgreen] (26,0) rectangle (28,2);
  \foreach \v in {20,22,25,28,35} \fill[sourceblue] (\v,0) circle (0.18);
  \node[font=\scriptsize] at (21.1,5.4) {36};
  \node[font=\scriptsize] at (23,3.4) {16};
  \node[font=\scriptsize,anchor=south] at (25.5,1) {1};
  \node[font=\scriptsize] at (27,1) {4};
  \node[font=\scriptsize] at (31.5,6) {81};
  \node[font=\scriptsize,anchor=west] at (29,2.5) {total area $=138$};
\end{tikzpicture}
\caption{Each squared deviation is the area of a square whose side is the distance
from a point to the mean (left squares red, right squares green; they overlap near
the mean). The variance is the total area, $138$, divided by $n-1=4$. The 35 ms
point alone contributes $81/138 \approx 59\%$ of it.}
\label{fig:center-spread-squares}
\end{figure}

The squares make one fact obvious: squaring weights large deviations heavily. A
point twice as far from the mean contributes four times as much. So the variance
and SD are \emph{not} robust: their breakdown point is $1/n$, like the mean. The
salary data have $s \approx 268$ thousand with the founder and $s \approx 7.8$
thousand without.

\begin{intuitionbox}
Variance is in squared units (ms$^2$, dollars$^2$), which no one can picture. The
square root returns to the original units. Read $s \approx 5.9$ ms as ``a typical
request is roughly 6 ms from the average.'' It is not exactly the average distance
(that would be the mean absolute deviation, $22/5 = 4.4$ ms here); squaring then
rooting gives extra weight to the far points.
\end{intuitionbox}

A useful shortcut for hand computation expands the square:
\[
  \sum (x_i - \xbar)^2 = \sum x_i^2 - n\xbar^2 .
\]
For the latencies, $\sum x_i^2 = 400+484+625+784+1225 = 3518$ and
$n\xbar^2 = 5\cdot 676 = 3380$, giving $138$ again. On a computer, avoid this
form: when $\xbar$ is large compared with $s$, the two big numbers nearly cancel and
floating-point error destroys the answer.

\section{Why divide by $n-1$?}
\label{sec:center-spread-bessel}

Dividing by $n$ would seem natural: the variance is ``the average squared
deviation.'' Indeed the \textbf{population variance} of a finite population of $N$
values with mean $\mu$ is
\[
  \sigma^2 = \frac{1}{N}\sum_{j=1}^N (x_j - \mu)^2 .
\]
Here the \textbf{population} is the whole collection we care about (every request
the server will ever handle), $\mu$ is its mean, and $\sigma^2$ is its average
squared distance from $\mu$. A \textbf{sample} is the handful of values we actually
see. The goal is to guess $\sigma^2$ from the sample.

The catch: we do not know $\mu$. We measure distances from $\xbar$ instead, and
$\xbar$ was computed from the same points we are measuring. It sits in the middle of
\emph{these} points, wherever they happen to fall, so they always look a little
closer to their center than they really are to the true center. Take the most
extreme case, two points. Their mean sits exactly halfway between them, so each is
half the gap away from it. The true mean $\mu$ is usually somewhere else, and then at
least one of the two points is more than half the gap from $\mu$. Distances from
$\xbar$ come out too small, every time, in the same direction. The rest of this
section measures exactly how much too small, and shows that the cure is to divide
by $n-1$.

\subsection{The sample mean hugs its own data}

Define the sum of squares about an arbitrary point $c$:
\[
  S(c) = \sum_{i=1}^n (x_i - c)^2 .
\]
Write $x_i - c = (x_i - \xbar) + (\xbar - c)$ and expand:
\[
  S(c) = \sum (x_i-\xbar)^2 + 2(\xbar - c)\underbrace{\sum (x_i - \xbar)}_{=\,0}
         + n(\xbar - c)^2
       = S(\xbar) + n(\xbar - c)^2 .
\]
So $S(c)$ is a parabola in $c$ with minimum exactly at $c = \xbar$. Whatever the true
mean $\mu$ is, $S(\mu) \ge S(\xbar)$, with a gap of $n(\xbar-\mu)^2$.

\begin{figure}[H]
\centering
\begin{tikzpicture}
\begin{axis}[
  width=10.5cm, height=5.8cm,
  domain=19:33, samples=100,
  xmin=19, xmax=33, ymin=60, ymax=420,
  axis lines=left, xlabel={center $c$ (ms)}, ylabel={$S(c)=\sum(x_i-c)^2$},
  tick label style={font=\scriptsize}, clip=false,
]
\addplot[distblue, thick] {138 + 5*(26-x)^2};
\addplot[only marks, mark=*, quotientorange] coordinates {(26,138)};
\addplot[only marks, mark=*, trapred] coordinates {(24,158)};
\draw[dashed,gray] (axis cs:19,138) -- (axis cs:33,138);
\draw[trapred,thick,{Stealth}-{Stealth}] (axis cs:24,140) -- (axis cs:24,156);
\node[trapred,font=\scriptsize,anchor=west] at (axis cs:19.3,112) {red: $S(\mu)=158$ if $\mu=24$; gap $n(\xbar-\mu)^2=20$};
\node[quotientorange,font=\scriptsize,anchor=west] at (axis cs:19.3,82) {orange: minimum $S(\xbar)=138$ at $c=\xbar=26$};
\end{axis}
\end{tikzpicture}
\caption{For the latency data, $S(c)=138+5(26-c)^2$. The sum of squares about
$\xbar$ is the smallest possible. If the true mean were $\mu=24$, the sum of squares
about $\mu$ would be $158$. Using $\xbar$ in place of $\mu$ can only make the sum
smaller, never larger.}
\label{fig:center-spread-hug}
\end{figure}

The identity has a plain reading. The squared distances to the true center equal
the squared distances to the sample's own center, \emph{plus} $n$ times the squared
distance between the two centers. That extra term is exactly what $\xbar$ hides.

\subsection{Every possible sample, listed}

Before any algebra, here is the bias in a case small enough to check by hand. A
population has three equally common values, $0$, $3$ and $6$. Its mean is $\mu=3$
and its variance is
\[
  \sigma^2 = \tfrac13\big((0-3)^2 + (3-3)^2 + (6-3)^2\big) = \tfrac13(9+0+9) = 6 .
\]
Draw a sample of $n=2$ values, each drawn independently (so the same value can come
up twice). There are $3\times 3 = 9$ equally likely ordered samples. For each one,
compute the sum of squares about the true $\mu=3$ and about the sample's own mean.

\begin{center}
\renewcommand{\arraystretch}{1.1}
\begin{tabular}{c c c c c}
\toprule
sample & $\xbar$ & $\sum (x_i-\mu)^2$ & $\sum (x_i-\xbar)^2$ & gap $2(\xbar-\mu)^2$ \\
\midrule
$0,0$ & $0$   & 18 & 0    & 18 \\
$0,3$ & $1.5$ & 9  & 4.5  & 4.5 \\
$0,6$ & $3$   & 18 & 18   & 0 \\
$3,0$ & $1.5$ & 9  & 4.5  & 4.5 \\
$3,3$ & $3$   & 0  & 0    & 0 \\
$3,6$ & $4.5$ & 9  & 4.5  & 4.5 \\
$6,0$ & $3$   & 18 & 18   & 0 \\
$6,3$ & $4.5$ & 9  & 4.5  & 4.5 \\
$6,6$ & $6$   & 18 & 0    & 18 \\
\midrule
average & $3$ & $12 = n\sigma^2$ & $6 = (n-1)\sigma^2$ & $6 = \sigma^2$ \\
\bottomrule
\end{tabular}
\end{center}

Read the bottom row. Measured from the true center, the sum of squares averages
$12 = 2\sigma^2$, so dividing by $n=2$ would be exactly right. Measured from $\xbar$,
it averages only $6$. In no row is it larger than the true-center column, and in
six of the nine rows it is smaller. Dividing by $n=2$ gives an average estimate of $3$, half
the truth. Dividing by $n-1=1$ gives $6$, exactly the truth.

\begin{figure}[H]
\centering
\begin{tikzpicture}[x=1.1cm,y=0.9cm]
  \foreach \yy/\lab in {2/{sample $0,6$},1/{sample $0,3$},0/{sample $0,0$}} {
    \draw[gray] (-0.5,\yy) -- (6.5,\yy);
    \node[anchor=east,font=\scriptsize] at (-0.6,\yy) {\lab};
  }
  \foreach \t in {0,3,6} \node[font=\scriptsize] at (\t,-0.55) {\t};
  \draw[dashed,structurepurple] (3,-0.3) -- (3,2.45);
  \node[structurepurple,font=\scriptsize,anchor=south] at (3,2.45) {true $\mu=3$};
  \fill[sourceblue] (0,2) circle (0.09);
  \fill[sourceblue] (6,2) circle (0.09);
  \fill[quotientorange] (3,2) -- ++(-0.1,-0.17) -- ++(0.2,0) -- cycle;
  \node[font=\scriptsize,anchor=west] at (6.7,2) {from $\xbar$: $18$; from $\mu$: $18$};
  \fill[sourceblue] (0,1) circle (0.09);
  \fill[sourceblue] (3,1) circle (0.09);
  \fill[quotientorange] (1.5,1) -- ++(-0.1,-0.17) -- ++(0.2,0) -- cycle;
  \node[font=\scriptsize,anchor=west] at (6.7,1) {from $\xbar$: $4.5$; from $\mu$: $9$};
  \fill[sourceblue] (0,0.12) circle (0.09);
  \fill[sourceblue] (0,-0.12) circle (0.09);
  \fill[quotientorange] (0.35,0) -- ++(-0.1,-0.17) -- ++(0.2,0) -- cycle;
  \node[font=\scriptsize,anchor=west] at (6.7,0) {from $\xbar$: $0$; from $\mu$: $18$};
\end{tikzpicture}
\caption{Three of the nine samples from the population $\{0,3,6\}$ (blue dots; the
orange triangle marks $\xbar$, drawn just right of the dots in the bottom row where
it coincides with them). When the sample happens to straddle $\mu$, nothing is lost.
When it lands to one side, $\xbar$ follows it and the spread seems to shrink, in the
bottom row all the way to zero.}
\label{fig:center-spread-enumerate}
\end{figure}

\subsection{The derivation}

Now do the same bookkeeping for any population and any $n$. Three words first.
An \textbf{estimator} is a recipe that turns a sample into a guess for an unknown
population number, such as ``divide the sum of squares by $n$.'' Because the sample
is random, the guess is random too. Its \textbf{expectation}, written $\E[\cdot]$, is
the long-run average of the guess over many repeated samples, exactly the
``average'' row of the table above. The \textbf{variance} of a random quantity,
$\Var(\cdot)$, is the expectation of its squared distance from its own expectation.

Treat the data as random: $X_1,\dots,X_n$ independent, each with mean $\mu$ and
variance $\sigma^2$ (capital letters mark values not yet drawn). Apply the identity
with $c=\mu$ and rearrange:
\[
  \sum_{i=1}^n (X_i - \Xbar)^2 = \sum_{i=1}^n (X_i - \mu)^2 - n(\Xbar - \mu)^2 .
\]
Take expectations of each term. (Expectation of a sum is the sum of expectations.)
\begin{itemize}
  \item $\E\big[(X_i - \mu)^2\big] = \sigma^2$ by definition of variance, so the first
  sum has expectation $n\sigma^2$.
  \item $\E\big[(\Xbar-\mu)^2\big] = \Var(\Xbar)$, because $\Xbar$ has expectation
  $\mu$. For independent values, variances add, so $\Var(X_1+\dots+X_n) = n\sigma^2$.
  Dividing a quantity by $n$ divides its variance by $n^2$, so
  $\Var(\Xbar) = n\sigma^2/n^2 = \sigma^2/n$. (Chapter~6 builds this result
  carefully; Probability From the Ground Up proves that variances of independent
  values add.) So the second term has expectation $n \cdot \sigma^2/n = \sigma^2$.
\end{itemize}
Therefore
\[
  \E\Big[\sum_{i=1}^n (X_i - \Xbar)^2\Big] = n\sigma^2 - \sigma^2 = (n-1)\sigma^2 .
\]
In the table: $n\sigma^2 = 12$, minus $\sigma^2 = 6$, leaves $(n-1)\sigma^2 = 6$.
Dividing the sum of squares by $n$ gives an estimator whose average is
$\frac{n-1}{n}\sigma^2$, too small. Dividing by $n-1$ gives
\[
  \E[s^2] = \sigma^2 .
\]
An estimator whose expectation equals the target is called \textbf{unbiased}; the
difference (expectation minus target) is its \textbf{bias}. The divide-by-$n$
estimator has bias $-\sigma^2/n$. The correction from $n$ to $n-1$ is
\textbf{Bessel's correction}.

\begin{historybox}
In the early 1800s astronomers combined many slightly disagreeing telescope readings
of a star's position and needed to report how precise the result was. Carl Friedrich
Gauss, in an 1823 treatise on combining observations, showed that the sum of squared
\emph{residuals} (each reading minus the fitted value) comes out too small, and that
dividing by the number of readings \emph{minus the number of quantities fitted from
them} removes the shortfall. For a single fitted mean that divisor is $n-1$. The name
honours his contemporary Friedrich Bessel, an astronomer who studied observation
errors. The same rule, ``subtract one per fitted parameter,'' gives the $n-2$ in
simple regression (Chapter~12).
\end{historybox}

\begin{whybox}
Two ways to see why exactly one $\sigma^2$ goes missing.

\emph{The wander of the center.} $\Xbar$ misses $\mu$ by a random amount whose
typical square is $\sigma^2/n$. Every one of the $n$ points is measured from the
wrong center, and the identity says the shortfall is $n(\Xbar-\mu)^2$. On average
that is $n\times\sigma^2/n = \sigma^2$: one data point's worth of variance.

\emph{Degrees of freedom.} The $n$ deviations $X_i - \Xbar$ must sum to zero, so
once you know any $n-1$ of them the last one is forced. With $n=2$ the two
deviations are always $+a$ and $-a$: one free number, not two. The sum of squares
contains only $n-1$ free pieces of information, called \textbf{degrees of freedom},
and we average over those. Each quantity estimated from the data before measuring
spread uses up one.
\end{whybox}

\subsection{Simulation}

Draw $n=3$ values from a Normal population with $\sigma^2 = 4$, compute both the
divide-by-$n$ estimate and $s^2$, and repeat. Figure~\ref{fig:center-spread-bessel-run}
tracks the running average of each estimator. The $n-1$ version settles near $4$;
the divide-by-$n$ version settles near $\tfrac{2}{3}\cdot 4 \approx 2.67$.

\begin{figure}[H]
\centering
\begin{tikzpicture}
\begin{axis}[
  width=12cm, height=6cm,
  xmode=log, xmin=1, xmax=2000, ymin=0.8, ymax=5,
  axis lines=left, xlabel={number of simulated samples (log scale)},
  ylabel={running average of estimate},
  legend pos=south east, legend style={font=\scriptsize},
  tick label style={font=\scriptsize},
]
\addplot[densityred, thick, mark=none] coordinates {(1,3.045) (2,2.476) (3,1.755) (5,1.195) (8,1.092) (10,1.748) (15,2.075) (20,2.166) (30,2.188) (40,2.414) (50,2.525) (70,2.750) (100,2.626) (130,2.513) (160,2.558) (200,2.603) (250,2.614) (300,2.648) (400,2.566) (500,2.595) (600,2.636) (700,2.586) (800,2.615) (1000,2.601) (1200,2.599) (1400,2.639) (1600,2.638) (1800,2.621) (2000,2.629)};
\addlegendentry{divide by $n$}
\addplot[distblue, thick, mark=none] coordinates {(1,4.568) (2,3.713) (3,2.633) (5,1.793) (8,1.639) (10,2.623) (15,3.112) (20,3.250) (30,3.282) (40,3.620) (50,3.787) (70,4.126) (100,3.939) (130,3.769) (160,3.836) (200,3.904) (250,3.921) (300,3.971) (400,3.849) (500,3.892) (600,3.954) (700,3.878) (800,3.923) (1000,3.902) (1200,3.899) (1400,3.959) (1600,3.957) (1800,3.931) (2000,3.943)};
\addlegendentry{divide by $n-1$ ($s^2$)}
\addplot[distblue, dashed, domain=1:2000] {4};
\addplot[densityred, dashed, domain=1:2000] {2.6667};
\node[font=\scriptsize,distblue,anchor=south] at (axis cs:400,4.25) {true $\sigma^2=4$};
\node[font=\scriptsize,densityred,anchor=north] at (axis cs:400,2.4) {$\tfrac{2}{3}\sigma^2\approx 2.67$};
\end{axis}
\end{tikzpicture}
\caption{2000 simulated samples of size $n=3$ from $\Normal(10, 4)$. Individual
estimates are noisy, but their long-run averages separate cleanly: $s^2$ averages
to the true variance, while dividing by $n$ averages to $\frac{n-1}{n}\sigma^2$.}
\label{fig:center-spread-bessel-run}
\end{figure}

\begin{figure}[H]
\centering
\begin{tikzpicture}
\begin{axis}[
  width=11cm, height=5.6cm,
  xmin=1.5, xmax=10.5, ymin=1.5, ymax=4.6,
  axis lines=left, xlabel={sample size $n$}, ylabel={average estimate},
  xtick={2,...,10},
  legend pos=south east, legend style={font=\scriptsize},
  tick label style={font=\scriptsize},
]
\addplot[only marks, mark=square*, distblue] coordinates {(2,3.979) (3,4.013) (4,4.016) (5,3.977) (6,4.022) (7,4.004) (8,3.999) (9,3.987) (10,4.003)};
\addlegendentry{$s^2$ (simulated)}
\addplot[only marks, mark=*, densityred] coordinates {(2,1.990) (3,2.675) (4,3.012) (5,3.182) (6,3.352) (7,3.432) (8,3.499) (9,3.544) (10,3.603)};
\addlegendentry{divide by $n$ (simulated)}
\addplot[densityred, thin, domain=2:10, samples=50] {4*(x-1)/x};
\addlegendentry{theory $4(n-1)/n$}
\addplot[distblue, thin, dashed, domain=2:10] {4};
\end{axis}
\end{tikzpicture}
\caption{Average of 50{,}000 estimates at each sample size. The divide-by-$n$
estimator follows the theoretical curve $\sigma^2(n-1)/n$: half the true variance at
$n=2$, still $10\%$ low at $n=10$. The bias vanishes as $n$ grows, which is why the
choice matters most for small samples.}
\label{fig:center-spread-bessel-n}
\end{figure}

\begin{trapbox}
Software disagrees on the default. NumPy's \texttt{np.var} and \texttt{np.std}
divide by $n$ unless you pass \texttt{ddof=1}; pandas' \texttt{.var()} and
\texttt{.std()} divide by $n-1$. The same data give different answers. (\texttt{ddof}
stands for ``delta degrees of freedom'': the divisor is $n-\texttt{ddof}$.) Also, $s$
itself is \emph{not} unbiased for $\sigma$: the square root bends downward, so
averaging and then rooting gives more than rooting and then averaging, and
$\E[s] < \sqrt{\E[s^2]} = \sigma$. For Normal data $s$ runs about $20\%$ low on
average at $n=2$, $6\%$ at $n=5$, and $3\%$ at $n=10$. Beyond small samples this is
usually ignored.
\end{trapbox}

When do you divide by $n$? Two cases. If the data \emph{are} the whole population
(every employee of the company, every request in the log you care about), there is
nothing larger to estimate, and the divide-by-$n$ formula is the population
variance. And if the true mean $\mu$ is known from outside the data, measure
deviations from $\mu$, not $\xbar$; then $\frac1n\sum(x_i-\mu)^2$ is unbiased (the
first column of the table). When the data are a sample used to learn about
something larger and $\mu$ is estimated by $\xbar$, divide by $n-1$.

\begin{recoverybox}
Pause and think: what does $s^2$ give for a sample of size $n=1$? (Zero over zero:
undefined.) Why is that the honest answer? (One value says nothing about spread;
dividing by $n$ would claim the spread is exactly $0$.)
\end{recoverybox}

\section{The key link: which center minimises which error?}
\label{sec:center-spread-loss}

The identity $S(c) = S(\xbar) + n(\xbar - c)^2$ says more than Bessel's correction.
It says:

\begin{theorem}[Mean minimises squared error]
For any data, $\displaystyle \xbar = \argmin_c \sum_{i=1}^n (x_i - c)^2$.
\end{theorem}

The median has the same property for a different loss.

\begin{theorem}[Median minimises absolute error]
For any data, a median minimises
\[
  A(c) = \sum_{i=1}^n \abs{x_i - c} .
\]
\end{theorem}

\begin{whybox}
Move $c$ a small step $\delta$ to the right. Every point to the right of $c$ gets
$\delta$ closer; every point to the left gets $\delta$ farther. So $A(c)$ changes by
$\delta \times (\#\text{left} - \#\text{right})$. While more points lie to the right,
moving right helps. The improvement stops exactly when the counts balance, that is,
at the median. Only the \emph{count} of points on each side matters, never their
distances, which is the same reason the median is robust.
\end{whybox}

Take the data $1, 2, 3, 4, 10$ (mean 4, median 3). Here $\sum x_i = 20$ and
$\sum x_i^2 = 130$, so
\[
  S(c) = 5c^2 - 40c + 130, \qquad S(4) = 50 .
\]
The absolute loss is piecewise linear: $A(0) = 20$, $A(1)=15$, $A(2)=12$, $A(3) = 11$,
$A(4) = 12$, $A(10) = 30$.

\begin{figure}[H]
\centering
\begin{tikzpicture}
\begin{axis}[
  name=left,
  width=6.6cm, height=5.4cm,
  xmin=0, xmax=10, ymin=0, ymax=240,
  axis lines=left, xlabel={$c$}, title={squared loss $S(c)$}, title style={font=\small},
  tick label style={font=\scriptsize},
]
\addplot[distblue, thick, domain=0:10, samples=80] {5*x^2-40*x+130};
\addplot[only marks, mark=*, quotientorange] coordinates {(4,50)};
\node[quotientorange,font=\scriptsize,anchor=north] at (axis cs:4,42) {min at mean $4$};
\addplot[only marks, mark=|, mark size=4pt, black] coordinates {(1,0)(2,0)(3,0)(4,0)(10,0)};
\end{axis}
\begin{axis}[
  at={(left.east)}, anchor=west, xshift=1.2cm,
  width=6.6cm, height=5.4cm,
  xmin=0, xmax=10, ymin=0, ymax=32,
  axis lines=left, xlabel={$c$}, title={absolute loss $A(c)$}, title style={font=\small},
  tick label style={font=\scriptsize},
]
\addplot[targetgreen, thick] coordinates {(0,20)(1,15)(2,12)(3,11)(4,12)(10,30)};
\addplot[only marks, mark=*, quotientorange] coordinates {(3,11)};
\node[quotientorange,font=\scriptsize,anchor=north] at (axis cs:3,9.5) {min at median $3$};
\addplot[only marks, mark=|, mark size=4pt, black] coordinates {(1,0)(2,0)(3,0)(4,0)(10,0)};
\end{axis}
\end{tikzpicture}
\caption{Total loss as a function of the summary $c$ for the data $1,2,3,4,10$
(ticks on the axis). Squared loss is a smooth parabola minimised at the mean; the
outlier at 10 pulls the minimum right. Absolute loss is piecewise linear with kinks
at the data and is minimised at the median.}
\label{fig:center-spread-loss}
\end{figure}

\begin{connectionbox}
This is the most important idea in the chapter for machine learning. Choosing a
loss function chooses which ``center'' your model predicts.
\begin{itemize}
  \item A regression model trained with \textbf{mean squared error} (MSE, $L_2$
  loss) learns to predict the conditional \emph{mean} of the target: the average
  label among inputs like the one in front of it. Linear
  regression (Chapter~12) is this idea with $c$ replaced by a line.
  \item A model trained with \textbf{mean absolute error} (MAE, $L_1$ loss) learns
  the conditional \emph{median}, and is less sensitive to outlying labels.
  \item \textbf{Quantile (pinball) loss} generalises $L_1$ to predict any quantile,
  such as p90 latency (Chapter~3 defines quantiles).
  \item \textbf{Huber loss} is quadratic near zero and linear far away: mean-like
  for small errors, median-like for outliers.
\end{itemize}
Picking the prediction that makes the average loss over the training data as small
as possible is called \emph{empirical risk minimisation}, and $\xbar$ is its
simplest instance: a model that predicts one constant for everyone. Chapter~15 returns to it.
\end{connectionbox}

\section{Other measures of spread: range, IQR, MAD}
\label{sec:center-spread-robust-spread}

The SD inherits the mean's sensitivity to outliers. Three alternatives trade
efficiency for robustness or simplicity. Take nine page-load times in ms, one of
which hit a cold cache:
\[
  12,\ 14,\ 15,\ 15,\ 17,\ 18,\ 20,\ 23,\ 90 .
\]
Mean $\approx 24.9$, median $17$, SD $s \approx 24.6$. Without the 90 the SD is
about $3.5$: one point inflated it sevenfold.

\begin{definition}[Range]
$\text{range} = x_{(n)} - x_{(1)}$.
\end{definition}
Here $90 - 12 = 78$. The range uses only the two most extreme values, so it has the
worst possible robustness and it grows with $n$ (more data, more chances for an
extreme). It is useful for sanity checks (``are all values physically
possible?''), not as a summary.

\begin{definition}[Interquartile range]
$\IQR = Q_3 - Q_1$, where $Q_1$ and $Q_3$ are the first and third quartiles: the
values below which $25\%$ and $75\%$ of the data fall.
\end{definition}
With nine sorted values the $25\%$ point is the 3rd value and the $75\%$ point is
the 7th (the default convention in NumPy and R; Chapter~3 compares conventions). So
$Q_1 = 15$, $Q_3 = 20$, and $\IQR = 5$. The IQR is the width of the middle half of
the data. Its breakdown point is $25\%$.

\begin{definition}[Median absolute deviation]
$\text{MAD} = \Med\big(\abs{x_i - \Med(x)}\big)$: the median distance from the median.
\end{definition}
The absolute distances from $17$ are $5, 3, 2, 2, 0, 1, 3, 6, 73$. Sorted:
$0,1,2,2,\mathbf{3},3,5,6,73$. The MAD is $3$. It is a median of a median, with
breakdown point $50\%$.

\begin{figure}[H]
\centering
\begin{tikzpicture}[x=0.13cm,y=0.8cm]
  \draw[thick] (8,0) -- (95,0);
  \foreach \t in {10,20,...,90} \draw (\t,0) -- (\t,-0.12) node[below,font=\scriptsize] {\t};
  \foreach \v/\h in {12/0.2,14/0.2,15/0.2,15/0.45,17/0.2,18/0.2,20/0.2,23/0.2,90/0.2}
    \fill[sourceblue] (\v,\h) circle (0.1cm);
  \draw[|-|,trapred,thick] (12,2.6) -- (90,2.6) node[midway,above,font=\scriptsize] {range $=78$};
  \draw[|-|,quotientorange,thick] (15,1.9) -- (20,1.9);
  \node[quotientorange,font=\scriptsize,anchor=west] at (22,1.9) {$\IQR = 20-15 = 5$};
  \draw[|-|,targetgreen,thick] (14,1.2) -- (20,1.2);
  \node[targetgreen,font=\scriptsize,anchor=west] at (22,1.2) {median $\pm$ MAD $= 17 \pm 3$};
  \draw[|-|,structurepurple,thick] (0.3,-1.1) -- (49.5,-1.1);
  \node[structurepurple,font=\scriptsize,anchor=west] at (51,-1.1) {mean $\pm$ SD $= 24.9 \pm 24.6$};
  \draw[dashed,gray] (17,0) -- (17,2.2);
\end{tikzpicture}
\caption{Four spread measures for the same nine values. The range and the SD band
are dominated by the single 90 ms value; the IQR and MAD describe where the bulk of
the data actually sits.}
\label{fig:center-spread-spread-measures}
\end{figure}

\begin{remark}
For Normal data the MAD is about $0.674\,\sigma$ and the IQR about $1.35\,\sigma$.
So $1.4826 \times \text{MAD}$ and $\IQR/1.35$ are robust estimates of $\sigma$.
Here $1.4826 \times 3 \approx 4.4$, close to the SD of the eight clean values
($3.5$) and nowhere near $24.6$. Robust outlier detectors use this: flag a point if
it is more than about $3 \times 1.4826 \times \text{MAD}$ from the median.
\end{remark}

\begin{lstlisting}[language=Python]
import numpy as np
from scipy import stats
x = np.array([12, 14, 15, 15, 17, 18, 20, 23, 90])
x.mean(), np.median(x)              # 24.89, 17.0
stats.trim_mean(x, 0.1)             # drops 0 from each end here (floor(0.9)=0)
x.std(ddof=1)                       # 24.64  (ddof=1 -> divide by n-1)
q1, q3 = np.percentile(x, [25, 75]) # 15.0, 20.0
stats.median_abs_deviation(x)       # 3.0
\end{lstlisting}

\section{Standardising: z-scores}
\label{sec:center-spread-z}

Alice scores 82 on an exam where the class mean is 70 and the SD is 8. Bob scores
90 on a different exam with mean 80 and SD 5. Who did better relative to their class?
Raw points cannot answer this: the exams have different centers and spreads.
Measure each score in units of its own SD from its own mean.

\begin{definition}[z-score]
The \textbf{z-score} of $x$ in a data set with mean $\xbar$ and SD $s$ is
\[
  z = \frac{x - \xbar}{s} .
\]
It says how many standard deviations $x$ lies above ($z>0$) or below ($z<0$) the mean.
\end{definition}

\[
  z_{\text{Alice}} = \frac{82 - 70}{8} = 1.5, \qquad
  z_{\text{Bob}} = \frac{90 - 80}{5} = 2.0 .
\]
Bob is further above his class.

\begin{recoverybox}
Pause and think: a third student scores 61 on Exam A. What is her z-score, and what
does its sign tell you? ($z = (61-70)/8 \approx -1.1$: about one SD below the mean.)
\end{recoverybox}

Standardising a whole data set (replacing every
$x_i$ by its $z_i$) produces values with mean $0$ and SD $1$, whatever the original
units; Section~\ref{sec:center-spread-linear} shows why.

\begin{figure}[H]
\centering
\begin{tikzpicture}
\begin{axis}[
  name=raw,
  width=12cm, height=4.4cm,
  xmin=40, xmax=105, ymin=0, ymax=0.095,
  axis lines=left, axis y line=none, xlabel={raw score},
  tick label style={font=\scriptsize}, clip=false,
  title={raw scores}, title style={font=\small},
]
\addplot[sourceblue, thick, domain=40:100, samples=120] {exp(-(x-70)^2/128)/(8*2.5066)};
\addplot[targetgreen, thick, domain=60:100, samples=120] {exp(-(x-80)^2/50)/(5*2.5066)};
\draw[sourceblue,thick] (axis cs:82,0) -- (axis cs:82,0.0162);
\draw[targetgreen,thick] (axis cs:90,0) -- (axis cs:90,0.0108);
\node[sourceblue,font=\scriptsize,anchor=east] at (axis cs:60,0.04) {Exam A: $70\pm 8$};
\node[targetgreen,font=\scriptsize,anchor=west] at (axis cs:86,0.075) {Exam B: $80\pm 5$};
\draw[sourceblue] (axis cs:82,0.0162) -- (axis cs:97,0.05) node[anchor=west,font=\scriptsize] {Alice 82};
\draw[targetgreen] (axis cs:90,0.0108) -- (axis cs:97,0.025) node[anchor=west,font=\scriptsize] {Bob 90};
\end{axis}
\begin{axis}[
  at={(raw.south)}, anchor=north, yshift=-2.3cm,
  width=12cm, height=4.4cm,
  xmin=-3.5, xmax=3.5, ymin=0, ymax=0.45,
  axis lines=left, axis y line=none, xlabel={z-score},
  tick label style={font=\scriptsize}, clip=false,
  title={after standardising, both classes lie on one scale}, title style={font=\small},
]
\addplot[black, thick, domain=-3.5:3.5, samples=120] {exp(-x^2/2)/2.5066};
\draw[sourceblue,thick] (axis cs:1.5,0) -- (axis cs:1.5,0.1295);
\draw[targetgreen,thick] (axis cs:2,0) -- (axis cs:2,0.054);
\draw[sourceblue] (axis cs:1.5,0.1295) -- (axis cs:2.3,0.30) node[anchor=west,font=\scriptsize] {Alice $z=1.5$};
\draw[targetgreen] (axis cs:2,0.054) -- (axis cs:2.6,0.17) node[anchor=west,font=\scriptsize] {Bob $z=2$};
\end{axis}
\end{tikzpicture}
\caption{Top: two exams with different means and spreads. Bob's 90 and Alice's 82
cannot be compared directly. Bottom: in SD units from each class mean, Bob ($z=2$)
is further into the upper tail than Alice ($z=1.5$).}
\label{fig:center-spread-z}
\end{figure}

\begin{connectionbox}
Feature standardisation in ML (\texttt{StandardScaler} in scikit-learn) is
exactly this: subtract each feature's mean, divide by its SD. Features measured in
metres and in milligrams end up on the same scale, so gradient descent converges
faster and size penalties on coefficients treat all features equally. Fit
the mean and SD on the training set only, then apply the same numbers to test data.
\end{connectionbox}

\section{The 68--95--99.7 rule}
\label{sec:center-spread-empirical}

What does ``$z = 2$'' mean in terms of how unusual a value is? It depends on the
shape. The most common shape is the \textbf{Normal} (or Gaussian) curve: a
symmetric bell, highest at the mean and falling off quickly on both sides. Heights,
measurement errors, and averages of many small effects tend to look like this
(Chapter~6 explains why; Probability From the Ground Up derives it). For data whose
histogram is roughly bell-shaped, the answer is fixed:

\begin{itemize}
  \item about $68\%$ of values lie within $1$ SD of the mean,
  \item about $95\%$ lie within $2$ SD,
  \item about $99.7\%$ lie within $3$ SD.
\end{itemize}
The exact Normal values are $68.27\%$, $95.45\%$, and $99.73\%$.

\begin{figure}[H]
\centering
\begin{tikzpicture}
\begin{axis}[
  width=12.2cm, height=6.4cm,
  xmin=-4, xmax=4, ymin=0, ymax=0.62,
  axis lines=left, axis y line=none,
  xtick={-3,-2,-1,0,1,2,3},
  xticklabels={$\mu-3\sigma$,$\mu-2\sigma$,$\mu-\sigma$,$\mu$,$\mu+\sigma$,$\mu+2\sigma$,$\mu+3\sigma$},
  tick label style={font=\scriptsize}, clip=false,
]
\addplot[fill=histfill!25, draw=none, domain=-3:3, samples=100] {exp(-x^2/2)/2.5066} \closedcycle;
\addplot[fill=histfill!55, draw=none, domain=-2:2, samples=100] {exp(-x^2/2)/2.5066} \closedcycle;
\addplot[fill=histfill, draw=none, domain=-1:1, samples=100] {exp(-x^2/2)/2.5066} \closedcycle;
\addplot[distblue, thick, domain=-4:4, samples=140] {exp(-x^2/2)/2.5066};
\draw[{Stealth}-{Stealth}] (axis cs:-1,0.45) -- (axis cs:1,0.45) node[midway,above,font=\scriptsize] {$68\%$};
\draw[{Stealth}-{Stealth}] (axis cs:-2,0.51) -- (axis cs:2,0.51) node[midway,above,font=\scriptsize] {$95\%$};
\draw[{Stealth}-{Stealth}] (axis cs:-3,0.57) -- (axis cs:3,0.57) node[midway,above,font=\scriptsize] {$99.7\%$};
\draw[gray,dotted] (axis cs:-3,0) -- (axis cs:-3,0.57);
\draw[gray,dotted] (axis cs:-2,0) -- (axis cs:-2,0.57);
\draw[gray,dotted] (axis cs:-1,0) -- (axis cs:-1,0.57);
\draw[gray,dotted] (axis cs:1,0) -- (axis cs:1,0.57);
\draw[gray,dotted] (axis cs:2,0) -- (axis cs:2,0.57);
\draw[gray,dotted] (axis cs:3,0) -- (axis cs:3,0.57);
\node[font=\scriptsize] at (axis cs:0,0.15) {$34.1\%+34.1\%$};
\node[font=\scriptsize] at (axis cs:1.5,0.04) {$13.6\%$};
\node[font=\scriptsize] at (axis cs:-1.5,0.04) {$13.6\%$};
\node[font=\scriptsize,anchor=south] at (axis cs:2.5,0.025) {$2.1\%$};
\node[font=\scriptsize,anchor=south] at (axis cs:-2.5,0.025) {$2.1\%$};
\end{axis}
\end{tikzpicture}
\caption{The 68--95--99.7 rule for a Normal distribution, with the exact share in
each band. Darker bands are closer to the mean. Each tail beyond $2\sigma$ holds
about $2.3\%$; the cut-off that leaves exactly $2.5\%$ in each tail is $1.96\sigma$,
the ``$1.96$'' in 95\% confidence intervals (Chapter~7).}
\label{fig:center-spread-6895997}
\end{figure}

\begin{example}
Adult heights in a population are roughly Normal with mean 170 cm and SD 7 cm.
Then about $68\%$ of adults are between 163 and 177 cm, about $95\%$ between 156 and
184 cm, and a height of 191 cm ($z=3$) is rarer than $0.15\%$ (one tail of the
$0.3\%$ outside $\pm 3$ SD).
\end{example}

\begin{trapbox}
The rule is for bell-shaped data only. For skewed data such as latencies it can be
badly wrong: a ``3-sigma'' alert threshold on a heavy-tailed metric fires far more
often than $0.3\%$ of the time. Plot the histogram before you use the rule.
\end{trapbox}

\section{Chebyshev's inequality: the guarantee for any data}
\label{sec:center-spread-chebyshev}

If you do not know the shape, is there anything you can say? Yes, and it needs only
the mean and SD.

\begin{theorem}[Chebyshev's inequality]
For any data set (or any distribution with finite variance) and any $k > 1$, the
fraction of values more than $k$ standard deviations from the mean is at most
$1/k^2$. Equivalently, at least $1 - 1/k^2$ lie within $k$ SDs.
\end{theorem}

\begin{whybox}
Let $p$ be the fraction of the $n$ values with $\abs{x_i - \mu} \ge k\sigma$ (using
the divide-by-$n$ $\sigma$ of the data). Each such value contributes at least
$(k\sigma)^2$ to the sum of squares, and the others contribute at least $0$. So
\[
  n\sigma^2 = \sum_{i} (x_i - \mu)^2 \ \ge\ (pn)(k\sigma)^2
  \quad\Longrightarrow\quad p \le \frac{1}{k^2}.
\]
The variance is a budget of squared distance. Far-away points are expensive, so
only a few can afford to be far.
\end{whybox}

\begin{figure}[H]
\centering
\begin{tikzpicture}
\begin{axis}[
  width=11.5cm, height=5.8cm,
  ybar, bar width=0.5cm,
  ymin=0, ymax=112,
  xtick={1,2,3,4}, xticklabels={1.5,2,3,4},
  xlabel={$k$ (number of SDs)}, ylabel={\% within $k$ SDs},
  enlarge x limits=0.18,
  legend style={at={(0.5,1.03)},anchor=south,legend columns=2,font=\scriptsize},
  nodes near coords, nodes near coords style={font=\tiny},
  tick label style={font=\scriptsize},
]
\addplot[fill=trapred!60, draw=trapred] coordinates {(1,55.6) (2,75.0) (3,88.9) (4,93.8)};
\addlegendentry{Chebyshev guarantee (any shape)}
\addplot[fill=histfill, draw=distblue] coordinates {(1,86.6) (2,95.4) (3,99.7) (4,99.99)};
\addlegendentry{Normal}
\end{axis}
\end{tikzpicture}
\caption{Minimum share of data within $k$ SDs of the mean. Chebyshev's bound holds
for every distribution, so it is conservative; Normal data do much better. At
$k=2$ Chebyshev promises at least $75\%$, the Normal gives $95\%$.}
\label{fig:center-spread-chebyshev}
\end{figure}

For the nine page-load times, $\xbar \approx 24.9$ and the population-style SD is
$\sigma \approx 23.2$, so the point $90$ is at $z \approx 2.8$. Chebyshev says at most
$1/2.8^2 \approx 13\%$ of points can be that far out; one of nine ($11\%$) is.
Chebyshev is tight only for contrived distributions, but it never fails. That is the
property you want when you know nothing about the shape. It also gives a one-line
proof of the law of large numbers (Chapter~6), the fact that a sample mean settles
down to the population mean as $n$ grows.

\section{Linear transformations of data}
\label{sec:center-spread-linear}

Unit conversions, currency conversions, and standardisation all have the form
$y_i = a + b x_i$. How do the summaries change?

\begin{theorem}[Shift and scale]
If $y_i = a + b\,x_i$ for all $i$, then
\[
  \ybar = a + b\,\xbar, \qquad s_y^2 = b^2 s_x^2, \qquad s_y = \abs{b}\, s_x .
\]
The median transforms like the mean, and the IQR, MAD, and range like the SD.
\end{theorem}

\begin{whybox}
Averaging is linear: $\frac1n\sum(a+bx_i) = a + b\xbar$. For the spread, look at the
deviations: $y_i - \ybar = (a+bx_i) - (a + b\xbar) = b(x_i - \xbar)$. The shift
$a$ cancels completely, because adding a constant moves every point and the center
together. The scale $b$ multiplies every deviation, so it multiplies the squared
deviations by $b^2$ and the SD by $\abs{b}$.
\end{whybox}

\begin{example}[Celsius to Fahrenheit]
Five daily highs in $^\circ$C: $18, 20, 21, 23, 28$. Mean $22^\circ$C, SD
$s \approx 3.81^\circ$C. With $F = 32 + 1.8\,C$ the readings become
$64.4, 68, 69.8, 73.4, 82.4$. Without recomputing,
\[
  \bar F = 32 + 1.8 \times 22 = 71.6^\circ\text{F}, \qquad
  s_F = 1.8 \times 3.81 \approx 6.85^\circ\text{F}.
\]
The $32$ does not touch the SD: a difference of temperatures has no offset.
\end{example}

\begin{figure}[H]
\centering
\begin{tikzpicture}
\begin{axis}[
  width=12cm, height=5.6cm,
  xmin=15, xmax=75, ymin=0.4, ymax=3.7,
  axis x line=bottom, axis y line=none,
  xlabel={value}, tick label style={font=\scriptsize}, clip=false,
]
\addplot[only marks, mark=*, sourceblue] coordinates {(20,3)(22,3)(25,3)(28,3)(35,3)};
\draw[|-|,sourceblue,thick] (axis cs:20.13,3.3) -- (axis cs:31.87,3.3);
\node[font=\scriptsize,anchor=west] at (axis cs:37,3) {$x$: mean 26, SD 5.87};
\addplot[only marks, mark=*, targetgreen] coordinates {(30,2)(32,2)(35,2)(38,2)(45,2)};
\draw[|-|,targetgreen,thick] (axis cs:30.13,2.3) -- (axis cs:41.87,2.3);
\node[font=\scriptsize,anchor=west] at (axis cs:47,2) {$x+10$: mean 36, SD 5.87};
\addplot[only marks, mark=*, quotientorange] coordinates {(40,1)(44,1)(50,1)(56,1)(70,1)};
\draw[|-|,quotientorange,thick] (axis cs:40.25,1.3) -- (axis cs:63.75,1.3);
\node[font=\scriptsize,anchor=east] at (axis cs:38,1) {$2x$: mean 52, SD 11.75};
\end{axis}
\end{tikzpicture}
\caption{The latency data, shifted and scaled. Bars show mean $\pm 1$ SD. Adding 10
slides everything right and leaves the spread unchanged. Doubling doubles both the
mean and the SD.}
\label{fig:center-spread-linear}
\end{figure}

Standardisation is the special case $a = -\xbar/s$, $b = 1/s$:
$z_i = (x_i - \xbar)/s$. The theorem gives $\bar z = -\xbar/s + \xbar/s = 0$ and
$s_z = s/s = 1$.

\begin{trapbox}
The shift-and-scale rule is for \emph{linear} maps only. For a nonlinear
transformation such as $\log$, the mean of the logs is not the log of the mean.
(The median does commute with any increasing function: the median of $\log x$ is
$\log$ of the median of $x$, because ranks are preserved.)
\end{trapbox}

\section{The mean of grouped data}
\label{sec:center-spread-grouped}

Dashboards and reports often give only a frequency table, not the raw values. Fifty
API response times were binned:

\begin{center}
\begin{tabular}{c c c r}
\toprule
bin (ms) & midpoint $m_j$ & count $f_j$ & $f_j m_j$ \\
\midrule
0--100   & 50  & 12 & 600 \\
100--200 & 150 & 25 & 3750 \\
200--300 & 250 & 8  & 2000 \\
300--500 & 400 & 5  & 2000 \\
\midrule
total    &     & 50 & 8350 \\
\bottomrule
\end{tabular}
\end{center}

Pretend every value in a bin sits at the bin's midpoint. Then
\[
  \xbar \approx \frac{\sum_j f_j m_j}{\sum_j f_j} = \frac{8350}{50} = 167 \text{ ms}.
\]
This is a \textbf{weighted mean}: each midpoint is weighted by how many observations
it represents. The same idea gives a grouped variance,
\[
  s^2 \approx \frac{1}{n-1}\sum_j f_j (m_j - \xbar)^2 = \frac{498{,}050}{49}
  \approx 10{,}164, \qquad s \approx 101 \text{ ms}.
\]

\begin{figure}[H]
\centering
\begin{tikzpicture}
\begin{axis}[
  width=11.5cm, height=5.6cm,
  xmin=0, xmax=520, ymin=0, ymax=30,
  axis lines=left, xlabel={response time (ms)},
  ylabel={count per 100 ms},
  xtick={0,100,200,300,400,500},
  tick label style={font=\scriptsize}, clip=false,
]
\addplot[ybar interval, fill=histfill, draw=distblue] coordinates {(0,12) (100,25) (200,8) (300,2.5) (500,2.5)};
\addplot[only marks, mark=*, mark size=1.8pt, black] coordinates {(50,12)(150,25)(250,8)(400,2.5)};
\draw[quotientorange, very thick] (axis cs:167,0) -- (axis cs:167,28);
\node[quotientorange, font=\scriptsize, anchor=west] at (axis cs:170,27) {grouped mean $\approx 167$};
\node[font=\scriptsize, anchor=south] at (axis cs:400,3) {5 values in 300--500 ms};
\end{axis}
\end{tikzpicture}
\caption{Histogram of the grouped response times. The last bin is twice as wide,
so its height is the count per 100 ms ($5/2 = 2.5$) to keep area proportional to
count. Dots mark the bin midpoints used to approximate the mean.}
\label{fig:center-spread-grouped}
\end{figure}

The approximation is exact if values are spread symmetrically within each bin, and
off otherwise. In the long right tail of skewed data the histogram falls across each
bin, so values in a bin cluster near its left edge and the midpoint rule tends to
overestimate. The grouped median can be read off by
finding the bin containing the $25$th and $26$th values: the 100--200 bin (values 13
through 37).

\begin{connectionbox}
The weighted mean $\sum_j w_j m_j$ with $w_j = f_j / n$ is an expectation in
disguise: the $w_j$ are probabilities that sum to $1$, and the formula is
$\E[X] = \sum_x x\, p(x)$. When an ML paper writes $\E_{x\sim p}[f(x)]$, read it as
``the $p$-weighted average of $f$.'' The sample mean is the special case where each
of the $n$ data points gets weight $1/n$.
\end{connectionbox}

\section{Choosing a summary}
\label{sec:center-spread-choosing}

\begin{figure}[H]
\centering
\begin{tikzpicture}[
  node distance=0.9cm and 0.6cm,
  q/.style={draw, rounded corners, fill=sourceblue!10, align=center, font=\small, text width=3.3cm},
  a/.style={draw, fill=targetgreen!12, align=center, font=\small, text width=3.4cm},
  >={Stealth}]
  \node[q] (cat) {Categorical data?};
  \node[a, right=1.2cm of cat] (mode) {mode\\(and counts)};
  \node[q, below=of cat] (out) {Outliers or strong skew?};
  \node[a, right=1.2cm of out] (rob) {median with IQR or MAD};
  \node[q, below=of out] (tot) {Need totals or a model with squared error?};
  \node[a, right=1.2cm of tot] (mean) {mean with SD};
  \node[a, below=of tot] (both) {report both mean and median; look at the plot};
  \draw[->] (cat) -- node[above,font=\scriptsize] {yes} (mode);
  \draw[->] (cat) -- node[left,font=\scriptsize] {no} (out);
  \draw[->] (out) -- node[above,font=\scriptsize] {yes} (rob);
  \draw[->] (out) -- node[left,font=\scriptsize] {no} (tot);
  \draw[->] (tot) -- node[above,font=\scriptsize] {yes} (mean);
  \draw[->] (tot) -- node[left,font=\scriptsize] {unsure} (both);
\end{tikzpicture}
\caption{A rough guide for picking a center and spread. The pairing matters: the SD
measures spread around the mean, and the IQR or MAD measure spread around the
median.}
\label{fig:center-spread-choose}
\end{figure}

When in doubt, compute both the mean and the median. If they are close, the data
are roughly symmetric and either will do. If they differ a lot, that difference is
itself the finding: the distribution is skewed or has outliers, and a single number
will mislead someone. Chapter~3 goes beyond a single center and spread: quantiles
such as p90 and p99, box plots that draw the IQR, and measures of shape.

\section{Summary}
\label{sec:center-spread-summary}

\begin{center}
\renewcommand{\arraystretch}{1.25}
\begin{tabular}{l l l l}
\toprule
measure & formula & minimises / meaning & breakdown \\
\midrule
mean $\xbar$ & $\frac1n\sum x_i$ & $\sum (x_i-c)^2$ & $1/n$ \\
median & middle sorted value & $\sum \abs{x_i-c}$ & $50\%$ \\
$\alpha$-trimmed mean & mean of middle & between the two & $\alpha$ \\
mode & most frequent & works for categories & -- \\
\midrule
variance $s^2$ & $\frac{1}{n-1}\sum (x_i-\xbar)^2$ & unbiased for $\sigma^2$ & $1/n$ \\
SD $s$ & $\sqrt{s^2}$ & typical distance, original units & $1/n$ \\
IQR & $Q_3 - Q_1$ & width of middle half & $25\%$ \\
MAD & $\Med\abs{x_i - \Med(x)}$ & robust spread & $50\%$ \\
range & $x_{(n)} - x_{(1)}$ & sanity check only & $1/n$ \\
\bottomrule
\end{tabular}
\end{center}

\begin{itemize}
  \item Deviations from the mean sum to zero: the mean is the balance point.
  \item Skew pulls the mean toward the long tail; for right skew,
  mode $<$ median $<$ mean.
  \item $\E\big[\sum(X_i-\Xbar)^2\big] = (n-1)\sigma^2$, because $\Xbar$ is fitted to
  the data. Hence divide by $n-1$.
  \item Squared loss leads to the mean; absolute loss leads to the median. Your choice
  of ML loss function chooses which one your model predicts.
  \item $z = (x-\xbar)/s$ measures distance in SD units. Bell-shaped data:
  $68/95/99.7\%$ within $1/2/3$ SDs. Any data: at least $1-1/k^2$ within $k$ SDs.
  \item $y = a + bx$ gives $\ybar = a + b\xbar$ and $s_y = \abs{b}s_x$.
  \item Grouped mean: $\sum f_j m_j / \sum f_j$, a weighted average of midpoints.
\end{itemize}

\section{Exercises}
\label{sec:center-spread-exercises}

\begin{enumerate}
  \item For the data $3, 5, 7, 7, 8$, compute the mean, median, mode, sample
  variance, and SD by hand using a deviations table.
  \item A data set has $n=10$ values. What is the smallest number of values you must
  replace to make the median arbitrarily large? What about the mean?
  \item Prove that $\sum_{i=1}^n (x_i - \xbar) = 0$ for any data.
  \item A city's daily highs over a month have mean $15^\circ$C and SD $4^\circ$C.
  Give the mean, SD, and variance in $^\circ$F ($F = 32 + 1.8C$).
  \item An IQ-style test has mean 100 and SD 15, and scores are roughly Normal. What
  is the z-score of 130? Using the 68--95--99.7 rule, roughly what share of people
  score above 130?
  \item Without knowing the shape of a distribution, what is the minimum fraction of
  values within $2.5$ SDs of the mean?
  \item With $n=2$ observations $X_1, X_2$ (independent, variance $\sigma^2$), show
  that $s^2 = (X_1 - X_2)^2/2$, and verify directly that $\E[s^2] = \sigma^2$.
  \item Estimate the mean from the grouped data: bin 0--10 has 4 values, 10--20 has
  10 values, 20--40 has 6 values.
  \item For the data $1, 2, 3, 4$, which values of $c$ minimise $\sum \abs{x_i - c}$?
  Which value minimises $\sum (x_i - c)^2$?
  \item Compute the MAD of $2, 4, 4, 5, 6, 9, 30$. Why is it so much smaller than
  the SD?
\end{enumerate}

\subsection*{Answers}

\begin{enumerate}
  \item Mean $30/5 = 6$; median $7$; mode $7$. Deviations $-3,-1,1,1,2$, squares
  $9,1,1,1,4$, sum $16$. $s^2 = 16/4 = 4$, $s = 2$.
  \item Five (half). With $5$ huge values they occupy sorted positions $6$--$10$, so
  the 6th value, and hence the median, is huge; with only $4$ the 5th and 6th values
  are untouched. For the mean, one value suffices.
  \item $\sum(x_i - \xbar) = \sum x_i - n\xbar = n\xbar - n\xbar = 0$.
  \item Mean $32 + 1.8 \times 15 = 59^\circ$F; SD $1.8 \times 4 = 7.2^\circ$F;
  variance $7.2^2 = 51.84$ ($^\circ$F)$^2$.
  \item $z = (130-100)/15 = 2$. About $5\%$ lie beyond $\pm 2$ SD, half of that above:
  roughly $2.5\%$ (exact Normal value $2.3\%$).
  \item $1 - 1/2.5^2 = 1 - 0.16 = 0.84$, so at least $84\%$.
  \item $\Xbar = (X_1+X_2)/2$, so $X_1 - \Xbar = (X_1-X_2)/2 = -(X_2-\Xbar)$ and the
  sum of squares is $2\cdot (X_1-X_2)^2/4$; dividing by $n-1 = 1$ gives
  $(X_1-X_2)^2/2$. Since $X_1 - X_2$ has mean $0$ and variance $2\sigma^2$,
  $\E[(X_1-X_2)^2] = 2\sigma^2$ and $\E[s^2] = \sigma^2$.
  \item Midpoints $5, 15, 30$: $(4\cdot 5 + 10\cdot 15 + 6 \cdot 30)/20 = 350/20 = 17.5$.
  \item Every $c$ in $[2, 3]$ minimises the absolute loss (two points on each side,
  so the slope is $0$ there); the squared loss is minimised only at $\xbar = 2.5$.
  \item Median $5$; absolute deviations $3,1,1,0,1,4,25$; sorted $0,1,1,1,3,4,25$;
  MAD $= 1$. The SD squares the deviation of the outlier $30$, so that one point
  dominates it (here $s \approx 9.7$); the MAD only uses the middle deviation.
\end{enumerate}
